%==============================================================================
%  CAPÍTULO VII
%==============================================================================
\chapter{Procedimiento concursal de liquidación ordinario}
\label{cap:liquidacion-ordinaria}

\fichaCapitulo{Realización colectiva de activos empresariales.}{Liquidación voluntaria y
forzosa, causales del artículo 117, audiencia inicial y juicio de oposición, desasimiento,
incautación, verificación y régimen recursivo restrictivo.}

%==============================================================================
\bloqueTeorico
%==============================================================================

\subsection{La liquidación como \lat{ultima ratio} del derecho de insolvencia}

La liquidación es el reconocimiento de que no hay valor de empresa en marcha que preservar,
o de que, habiéndolo, los acreedores han decidido no preservarlo. Frente a la
reorganización ---que apuesta a la continuidad--- la liquidación acepta la desaparición de
la unidad económica y organiza lo único que resta: la conversión ordenada del activo en
dinero y su distribución conforme a las preferencias legales.

Esta función explica los tres rasgos que estructuran el procedimiento y que lo distinguen
del de reorganización:

\begin{enumerate}
  \item \textbf{Admite iniciativa del acreedor.} A diferencia de la reorganización, que
  sólo puede pedir el deudor, la liquidación puede ser forzosa: un acreedor puede
  demandarla cuando concurre una causal legal.

  \item \textbf{Produce desasimiento.} El deudor pierde la administración de sus bienes,
  que pasa de pleno derecho al liquidador. No es una intervención ni una vigilancia: es un
  desplazamiento total de la facultad de administrar.

  \item \textbf{Restringe los recursos.} El régimen recursivo es deliberadamente estrecho,
  al servicio de la celeridad de la realización, y ha sido por ello objeto de reproche
  constitucional (sección \ref{sec:cap07-const}).
\end{enumerate}

\subsection{Desasimiento y administración privativa del liquidador}

El desasimiento es la institución central del procedimiento y su efecto más intenso sobre
el deudor. Conviene precisar desde ya lo que es y lo que no es. No priva al deudor de la
propiedad de sus bienes: el fallido sigue siendo dueño. Lo priva de la \emph{facultad de
administrarlos y disponer de ellos}, que se traslada íntegramente al liquidador.

La consecuencia, regulada en el \art{130}, es drástica: los actos y contratos que el deudor
ejecute o celebre sobre los bienes sujetos al procedimiento después de la resolución son
\destacar{nulos}. No inoponibles ni anulables: nulos. La administración pasó a otro, y quien
ya no administra no puede obligar válidamente el patrimonio.

\begin{foco}[Lo que el cliente debe entender el primer día]
Desde la resolución de liquidación, el deudor no puede vender, pagar, cobrar ni gravar sus
bienes: cualquier acto que celebre es nulo. Explicarlo antes de la resolución evita que el
cliente, actuando de buena fe, celebre operaciones nulas que además pueden leerse como
distracción de bienes.
\end{foco}

%==============================================================================
\bloqueNormativo
%==============================================================================

\subsection{La liquidación voluntaria}

El \art{115} permite a la \empresadeudora{} solicitar su propia liquidación ante el juzgado
de letras competente, acompañando, con copia, un conjunto de antecedentes cuyo primero
concentra la información patrimonial:

\begin{enumerate}
  \item \textbf{Lista de bienes}, lugar en que se encuentran y gravámenes que los afectan,
  \destacar{incluyendo los que estén en su poder en calidad distinta de la de dueño} y los
  bienes constituidos en garantía a su favor, con la documentación que lo acredite. Debe
  indicar, además, su participación en sociedades, comunidades y comunidades hereditarias.
\end{enumerate}

\verificar{Completar el listado de antecedentes del artículo 115 con los numerales
restantes: lista de acreedores, memoria de las causas de la crisis, y demás documentación
según el régimen tributario del deudor.}

Presentada la solicitud, el \art{116} ordena al tribunal revisarla y, si cumple los
requisitos, proceder dentro de tercero día conforme a los \art{37} y \art{129}, esto es,
requerir la nominación del liquidador y dictar la Resolución de Liquidación.

\begin{practica}[Los bienes de terceros deben declararse desde el inicio]
El \art{115} exige declarar los bienes que el deudor tiene en calidad distinta de la de
dueño ---en depósito, comodato, \emph{leasing}, consignación---. Identificarlos en la
solicitud evita que se incauten como propios y previene las tercerías. La omisión no
beneficia al deudor: sólo genera conflicto con el tercero y desgaste del procedimiento.
\end{practica}

\subsection{La liquidación forzosa: las causales del artículo 117}

Cualquier acreedor puede demandar la liquidación de una \empresadeudora{} en tres casos
tasados:

\begin{description}[leftmargin=0pt,style=nextline,font=\sffamily\bfseries]
  \item[Causal 1 — Cesación en el pago de un título ejecutivo] Si el deudor cesa en el pago
  de una obligación que conste en \destacar{título ejecutivo vencido}, que sea una
  obligación \destacar{propia de la actividad} de la empresa deudora con el acreedor
  solicitante.

  La ley agrega una salvaguarda esencial: esta causal \destacar{no puede invocarse contra
  los fiadores, codeudores solidarios o subsidiarios, o avalistas} de la empresa deudora
  que cesó en el pago de las obligaciones garantizadas por ellos.

  \item[Causal 2 — Dos ejecuciones sin bienes suficientes] Si existen dos o más títulos
  ejecutivos vencidos, provenientes de \destacar{obligaciones diversas}, con al menos dos
  ejecuciones iniciadas, y el deudor no presentó bienes suficientes para responder a la
  prestación y sus costas dentro de los \destacar{cuatro días} siguientes a los respectivos
  requerimientos.

  \item[Causal 3 — Deudor no habido] Cuando la empresa deudora o sus administradores no
  sean habidos y hayan dejado cerradas sus oficinas o establecimientos sin haber nombrado
  mandatario con facultades suficientes para cumplir sus obligaciones y contestar demandas.
  En este caso el demandante puede invocar incluso un crédito sujeto a plazo o condición
  suspensiva.
\end{description}

\begin{foco}[La salvaguarda del fiador es un argumento de defensa]
La causal 1 no alcanza al fiador ni al codeudor. Con frecuencia el acreedor demanda la
liquidación del garante para presionar el pago; la defensa consiste en oponer que la
obligación garantizada es del deudor principal y que la causal expresamente excluye al
fiador. Es una excepción de texto, no de interpretación.
\end{foco}

\subsection{Requisitos de la demanda y la consignación del artículo 118}

El \art{118} exige que la demanda señale la causal y sus hechos justificativos, y acompañe:

\begin{enumerate}
  \item Los documentos que acrediten la causal invocada.
  \item \textbf{Vale vista o boleta bancaria} a la orden del tribunal por
  \destacar{100 unidades de fomento}, para subvenir los gastos iniciales y los honorarios
  del liquidador. Dictada la Resolución de Liquidación, esa suma se considera un crédito del
  acreedor solicitante contra el deudor.
\end{enumerate}

\begin{alerta}[El costo de entrada de la liquidación forzosa]
Demandar la liquidación cuesta 100 UF por adelantado. El acreedor que la usa como mecanismo
de presión debe saber que ese desembolso es real y que sólo recupera como crédito si la
liquidación efectivamente se declara. Es un filtro deliberado contra el uso abusivo del
procedimiento.
\end{alerta}

\subsection{La audiencia inicial y las opciones del deudor}

El \art{120} regula la audiencia inicial, que es el momento decisivo del juicio de
oposición. El tribunal informa al deudor de la demanda y de los efectos de una eventual
liquidación. Acto seguido, el deudor \destacar{debe señalar el nombre, domicilio y correo
de sus tres mayores acreedores según su contabilidad}; si no lo hace, el tribunal tiene por
no presentada la actuación y dicta de inmediato la Resolución de Liquidación.

Cumplida esa carga, el deudor puede optar por una de cuatro actuaciones:

\begin{description}[leftmargin=0pt,style=nextline,font=\sffamily\bfseries]
  \item[a) Consignar] Fondos suficientes para pagar el crédito demandado y las costas. El
  tribunal tiene por efectuada la consignación, ordena liquidar el crédito y tasar las
  costas, y fija plazo de pago. Si el deudor no paga en ese plazo, se dicta la Resolución de
  Liquidación.

  \item[b) Allanarse] Por escrito o verbalmente, caso en que el tribunal dicta la
  Resolución de Liquidación.

  \item[c) Acogerse a reorganización] Del Capítulo III (capítulo
  \ref{cap:reorganizacion-ordinaria}). Es la vía por la que una demanda de liquidación puede
  reconducirse a un salvataje.

  \item[d) Oponerse] En cuyo caso se abre el juicio de oposición conforme al Párrafo 3.
\end{description}

\begin{foco}[La decisión estratégica de la audiencia inicial]
Las cuatro opciones no son equivalentes ni reversibles. Consignar enerva la causal pero
supone tener los fondos. Allanarse ordena la salida cuando la liquidación es inevitable.
Acogerse a reorganización es la jugada de quien tiene un negocio viable y llega preparado.
Oponerse exige prueba de solvencia lista. La decisión debe estar tomada \emph{antes} de la
audiencia, porque en ella no hay tiempo de improvisar. Véase el capítulo
\ref{cap:tram-liq-ordinaria}.
\end{foco}

\subsection{La Resolución de Liquidación}

El \art{129} detalla el contenido de la Resolución de Liquidación, que además de las
menciones de los \art[CPC]{169} y \art[CPC]{170} debe contener:

\begin{enumerate}
  \item Las consideraciones de hecho y de derecho que fundan el rechazo de las excepciones
  del deudor, si es procedente.
  \item La determinación de que el deudor es una empresa deudora, individualizándola.
  \item La designación de liquidador titular y suplente, provisionales, y la
  \destacar{orden de incautar todos los bienes} del deudor.
\end{enumerate}

\verificar{Completar los numerales restantes del artículo 129, en particular la orden de
publicación y el llamado a junta constitutiva.}

\subsection{Efectos de la resolución: el desasimiento}

El \art{130} enumera los efectos sobre el deudor y sus bienes desde la dictación de la
resolución:

\begin{enumerate}
  \item \textbf{Inhibición de administrar} (\Nro 1). El deudor queda inhibido de pleno
  derecho de la administración de todos sus bienes presentes ---los existentes a la época de
  la resolución, excluidos los inembargables---. La administración pasa de pleno derecho al
  liquidador. \destacar{Son nulos} los actos y contratos posteriores del deudor sobre esos
  bienes.
  \item \textbf{Subsistencia de la propiedad} (\Nro 2). El desasimiento no priva al deudor
  del dominio, sino de la facultad de administrar y disponer.
\end{enumerate}

Los demás efectos de la apertura ---la fijación irrevocable de los derechos de los acreedores
(\art{134}) y la suspensión de las ejecuciones individuales (\art{135})--- se desarrollan más
adelante en este mismo capítulo.

\subsection{Incautación e inventario}

Asumido el cargo, el \art{163} ordena al liquidador adoptar de inmediato las medidas
conservativas necesarias y practicar la diligencia de incautación y confección del
inventario. El \art{164} regula el acta de incautación, que debe contener, entre otras
menciones, la singularización de los domicilios, la fecha y los asistentes, la constancia
de si fue necesario el auxilio de la fuerza pública, y \destacar{la constancia de todo
derecho o pretensión formulados por terceros} sobre los bienes.

\begin{practica}[La constancia de las pretensiones de terceros]
El acta debe registrar toda pretensión de tercero sobre un bien incautado. Esa constancia
es el punto de partida de las eventuales tercerías y la prueba de que el liquidador actuó
diligentemente. El abogado del deudor debe velar por que las pretensiones legítimas queden
consignadas en el acto, no después.
\end{practica}

\subsection{Verificación ordinaria de créditos}

El \art{170} concede a los acreedores un plazo de \destacar{treinta días} contado desde la
notificación de la Resolución de Liquidación para verificar sus créditos y alegar su
preferencia, acompañando los títulos justificativos e indicando correo electrónico para las
notificaciones. Vencido el plazo, dentro de los dos días siguientes el liquidador publica en
el \boletin{} la nómina correspondiente.

\begin{alerta}[Treinta días, no quince]
El plazo de verificación en la liquidación es de treinta días, distinto del de quince de la
reorganización (\art{70}). La confusión entre ambos es una fuente de verificaciones
extemporáneas.
\end{alerta}

\subsubsection{Los acreedores de servicios de utilidad pública}

El \art{171} contiene una regla protectora de la continuidad: los prestadores de servicios de
utilidad pública ---electricidad, agua, gas, telecomunicaciones--- deben verificar los créditos
por suministros anteriores a la resolución, y \destacar{no pueden suspender el servicio} con
posterioridad a ella salvo autorización del tribunal, previa audiencia del liquidador. Los
suministros posteriores a la resolución son de cargo de la masa.

\begin{practica}[El corte de suministro no procede sin autorización judicial]
Si tras la liquidación la empresa de servicios amenaza con cortar el suministro por deuda
anterior, la defensa es directa: el \art{171} se lo prohíbe sin autorización del tribunal. Es
una herramienta para preservar la operación durante la continuación de actividades.
\end{practica}

\subsubsection{Objeción e impugnación de créditos}

Verificados los créditos, se abre la fase de depuración. El \art{174} concede a los acreedores,
al liquidador y al deudor un plazo de \destacar{diez días} desde el vencimiento de la
verificación ordinaria para deducir \destacar{objeción fundada} sobre la existencia, el monto o
la preferencia de los créditos. Vencido ese plazo sin objeciones, los créditos se tienen por
reconocidos.

Formuladas objeciones, el \art{175} encarga al liquidador arbitrar el ajuste entre los
interesados. Si las objeciones no se subsanan, los créditos quedan \destacar{impugnados}: el
liquidador los acumula, emite un informe sobre si existen fundamentos plausibles y acompaña la
nómina al tribunal, que resuelve.

\begin{foco}[La objeción es del acreedor; la impugnación la construye el liquidador]
El sistema distingue dos fases. Primero, cualquier interesado \emph{objeta} en diez días.
Después, si la objeción no se subsana, el crédito pasa a \emph{impugnado} y decide el tribunal
con el informe del liquidador. Para el acreedor que quiere excluir un crédito inflado de un
competidor ---o para el deudor que quiere depurar su pasivo--- la ventana de los diez días del
\art{174} es fatal: dejarla pasar consolida el crédito objetable.
\end{foco}

\subsubsection{Verificación extraordinaria}

El \art{179} permite al acreedor que no verificó en el período ordinario hacerlo
\destacar{mientras no esté firme la cuenta final de administración}, pero sólo para ser
considerado en los \destacar{repartos futuros} y aceptando todo lo obrado con anterioridad. Los
créditos verificados extraordinariamente pueden objetarse e impugnarse conforme a los \art{174}
y \art{175}.

\begin{alerta}[Verificar tarde cuesta los repartos ya hechos]
El acreedor que verifica extraordinariamente no recupera lo repartido antes de su comparecencia:
sólo participa de lo que reste. La diligencia en verificar dentro del plazo ordinario tiene, por
tanto, consecuencia patrimonial directa.
\end{alerta}

\subsection{Fijación de derechos y suspensión de las ejecuciones individuales}

Dos efectos de la resolución completan el cuadro del \desasimiento{}. El \art{134} dispone que
la Resolución de Liquidación \destacar{fija irrevocablemente los derechos de todos los
acreedores en el estado que tenían al día de su pronunciamiento}, salvo las excepciones
legales: se congela la masa pasiva a esa fecha. Y el \art{135} \destacar{suspende el derecho de
los acreedores a ejecutar individualmente} al deudor ---con la excepción de los acreedores
hipotecarios y prendarios, que pueden deducir o continuar sus acciones sobre los bienes
gravados---.

\begin{foco}[La liquidación colectiviza el cobro]
Desde la resolución, ningún acreedor cobra por su cuenta: todos concurren al reparto colectivo.
La única salvedad son los acreedores con garantía real sobre bienes determinados, que conservan
su acción individual sobre esos bienes. Identificar quién tiene garantía real ---y sobre qué
bien--- es, por tanto, el primer trabajo de dimensionamiento de la masa.
\end{foco}

\subsection{Exigibilidad anticipada y valor actual de los créditos}

El \art{136} produce un efecto de gran alcance: dictada la resolución, \destacar{todas las
obligaciones dinerarias del deudor se entienden vencidas y actualmente exigibles}, para que los
acreedores puedan verificarlas y percibir el pago. Las obligaciones a plazo se anticipan; se
pagan según su valor actual más reajustes e intereses.

El \art{137} fija las reglas de cálculo del \destacar{valor actual} de los créditos ---
distinguiendo entre reajustables y no reajustables, en moneda nacional o extranjera, con o sin
intereses---, para traer todos los créditos a una fecha común y hacerlos comparables en el
reparto.

\begin{foco}[La liquidación acelera todo el pasivo]
El crédito a tres años vence hoy. Para el acreedor con plazo pendiente, la resolución de
liquidación adelanta su exigibilidad ---y lo obliga a verificar ya---; para el deudor, significa
que ningún plazo lo protege. Es una manifestación más de que la liquidación cristaliza la
situación patrimonial a la fecha de su apertura.
\end{foco}

\subsection{Compensación y obligaciones conexas}

El \art{140} restringe la compensación: la resolución \destacar{impide toda compensación que no
hubiere operado antes por el solo ministerio de la ley} entre las obligaciones recíprocas del
deudor y sus acreedores. La excepción son las \destacar{obligaciones conexas} ---derivadas de un
mismo contrato o negociación--- que sí pueden compensarse aunque sean exigibles en plazos
distintos.

\begin{practica}[La conexidad rescata la compensación]
El acreedor que es a su vez deudor de la masa pierde, en principio, la compensación tras la
apertura. Pero si ambas obligaciones nacen del mismo contrato ---por ejemplo, un suministro con
saldos recíprocos--- la conexidad la preserva. Alegar y probar la conexidad es la vía para no
tener que pagar íntegro a la masa y cobrar sólo un dividendo.
\end{practica}

\subsection{La acumulación de juicios}

El \art{142} consagra la regla general: \destacar{todos los juicios civiles pendientes contra el
deudor se acumulan} al procedimiento de liquidación, y los posteriores se promueven ante el
tribunal del concurso. El \art{143} exceptúa los juicios ante árbitros, los de arbitraje forzoso
y los sometidos a tribunales especiales. Los \art{144} a \art{146} regulan la acumulación de los
juicios ejecutivos, con reglas propias según haya o no excepciones opuestas y según existan
codemandados distintos del deudor.

\begin{foco}[El concurso atrae los litigios del deudor]
La apertura concentra ante un solo tribunal la conflictividad dispersa del deudor. Para el
acreedor que litigaba en otra sede, esto significa que su juicio se traslada ---salvo las
excepciones del art. 143--- al tribunal del concurso. Verificar el estado de cada juicio
pendiente y su suerte bajo la acumulación es parte del mapa inicial de la liquidación.
\end{foco}

\subsection{Reivindicación, retención y contratos}

La ley preserva los derechos de ciertos terceros sobre bienes en poder del deudor:

\begin{description}[leftmargin=0pt,style=nextline,font=\sffamily\bfseries]
  \item[Reivindicación de efectos de comercio (\art{151})] Pueden reivindicarse los efectos de
  comercio y documentos de crédito no pagados, existentes a la apertura, que el propietario
  entregó al deudor por un título no traslaticio de dominio.
  \item[Reivindicación de mercaderías (\art{152})] Procede respecto de mercaderías en
  situaciones análogas.
  \item[Derecho legal de retención (\art{160})] Tiene lugar cuando quien pagó o se obligó a
  pagar por el deudor tiene en su poder mercaderías o valores de éste, si la tenencia nace de un
  hecho voluntario anterior.
  \item[Arrendamiento (\art{141})] El derecho legal de retención del arrendador no puede
  declararse después de la resolución, y por treinta días el arrendador no puede realizar los
  bienes muebles destinados al giro por las rentas vencidas.
\end{description}

\begin{alerta}[Los bienes de terceros no integran la masa]
La reivindicación y la retención confirman lo ya señalado a propósito del inventario: no todo lo
que está en poder del deudor es suyo. Los efectos de comercio entregados en cobranza, las
mercaderías en consignación y los bienes retenidos por terceros pueden salir de la masa. La
identificación temprana de estos bienes ---y de sus titulares--- evita incautaciones indebidas y
tercerías.
\end{alerta}

\subsection{El deber de colaboración del deudor}

El \art{169} impone al deudor el deber de \destacar{indicar y poner a disposición del liquidador
todos sus bienes y antecedentes}, bajo apercibimiento de \destacar{arresto hasta por dos meses o
multa de hasta 10 UTM}. Es la contrapartida del desasimiento: el deudor pierde la administración,
pero queda obligado a colaborar con quien la asume.

\begin{practica}[El deber de colaboración y el incidente de mala fe]
La negativa a colaborar no sólo expone al arresto del \art{169}: alimenta el incidente de mala
fe del \art{169 A} (capítulo \ref{cap:mala-fe}). Al cliente deudor debe explicársele que la
colaboración diligente ---entrega de documentación, individualización de bienes--- es también su
mejor defensa frente a la denegación del \discharge{}.
\end{practica}

\subsection{Las juntas de acreedores}

El \art{180} organiza la deliberación colectiva en tres tipos de junta: \destacar{constitutiva,
ordinarias y extraordinarias}. El \art{181} fija el quórum para sesionar: toda junta se entiende
legalmente constituida con la concurrencia de acreedores que representen al menos el
\destacar{25\,\% del pasivo con derecho a voto}.

\subsubsection{La Junta Constitutiva}

Es la primera junta del procedimiento. El \art{193} la sitúa al \destacar{trigésimo segundo día}
contado desde la publicación de la Resolución de Liquidación en el \boletin{}. En ella los
acreedores ratifican o sustituyen al liquidador, deciden sobre la continuidad del giro y fijan
las modalidades de realización.

\subsubsection{La inasistencia y la ratificación de pleno derecho}

El \art{195} regula un supuesto de la mayor importancia práctica: si en la \emph{segunda
citación} no asiste ningún acreedor con derecho a voto, el secretario del tribunal lo certifica
y, \destacar{sin necesidad de declaración judicial}, los liquidadores provisionales quedan
ratificados de pleno derecho como definitivos.

\begin{foco}[La pasividad de los acreedores consolida al liquidador]
En los concursos de escaso activo ---la mayoría--- los acreedores rara vez concurren. La
consecuencia del \art{195} es que el liquidador propuesto por el peticionario tiende a
consolidarse por inasistencia. Para el acreedor que quiere incidir en quién administra, la única
oportunidad es comparecer a la junta constitutiva: después, el liquidador ya está ratificado.
\end{foco}

\subsubsection{Publicidad, acta y derecho a voz}

El \art{182} declara las juntas \destacar{públicas} y reconoce derecho a voz a todos los
acreedores que verificaron ---tengan o no derecho a voto---, al liquidador, al deudor y al
Superintendente. El \art{183} exige suscribir la nómina de asistencia; el \art{184} obliga a
levantar acta de todo lo obrado ---acuerdos adoptados y propuestas desestimadas--- y a
\destacar{publicarla en el \boletin{} al día siguiente}.

\subsubsection{La determinación del derecho a voto}

El \art{190} regula la audiencia en que el tribunal determina el \destacar{derecho a voto} de
los acreedores cuyos créditos no están reconocidos, celebrada el día anterior a la junta o el
mismo día antes de la constitutiva. Es un filtro decisivo: quién vota, y por cuánto, se decide
aquí, y de ello dependen todos los quórums.

\begin{practica}[El derecho a voto se pelea antes de la junta]
Las mayorías se ganan o se pierden en la audiencia de determinación del derecho a voto, no en la
junta misma. El acreedor que quiere pesar ---o excluir a un competidor de voto dudoso--- debe
comparecer a esa audiencia con sus objeciones preparadas. Llegar a la junta sin haber litigado el
derecho a voto es llegar tarde.
\end{practica}

\subsubsection{Materias de cada tipo de junta}

\begin{description}[leftmargin=0pt,style=nextline,font=\sffamily\bfseries]
  \item[Junta Constitutiva (\art{196}-\art{197})] Presidida por el juez. El liquidador presenta
  una cuenta escrita y verbal del estado de los negocios, el activo y el pasivo, y la gestión
  realizada. Se ratifica o sustituye al liquidador y se fijan las bases de la realización.
  \item[Primera Junta Ordinaria (\art{198})] Trata, si no se acordaron antes, el informe
  actualizado de activo y pasivo y el \destacar{plan de realización} de los bienes.
  \item[Junta Extraordinaria (\art{199}-\art{200})] Procede por orden del tribunal, a petición
  del liquidador o de la \superir{}, o \destacar{a solicitud de acreedores que representen al
  menos el 25\,\% del pasivo con derecho a voto}. Son materias exclusivas suyas la
  \destacar{revocación de los liquidadores} definitivos y la presentación de propuestas de
  acuerdo de reorganización (el rescate del \art{257}).
\end{description}

\begin{foco}[El 25\,\% que convoca la extraordinaria]
Un acreedor o grupo que reúna el 25\,\% del pasivo con derecho a voto puede forzar una junta
extraordinaria ---y en ella, plantear la revocación del liquidador o un acuerdo de
reorganización que rescate la empresa---. Es la palanca de los acreedores relevantes frente a un
liquidador que no los representa o a una liquidación que preferirían reconducir.
\end{foco}

\subsection{Pago administrativo y repartos de fondos}

\subsubsection{Pago administrativo de los créditos preferentes}

El \art{244} permite al liquidador \destacar{pagar administrativamente} ---tan pronto haya
fondos suficientes y asegurados los gastos y los créditos de mejor derecho--- los créditos del
\artcc{2472}. Los de los números 1 y 4 (costas y ciertos créditos laborales) pueden pagarse
\destacar{sin necesidad de verificación}, lo que agiliza la satisfacción del crédito laboral
preferente (capítulo \ref{cap:creditos-laborales}).

\subsubsection{La propuesta de reparto}

El \art{247} sujeta la propuesta de reparto a requisitos copulativos: disponibilidad de fondos
para abonar a los acreedores reconocidos \destacar{no menos del 5\,\% de sus acreencias} ---salvo
acuerdo de junta por un porcentaje inferior--- y reserva previa de los dineros para los gastos
del procedimiento.

\subsubsection{Compensación de créditos recíprocos}

El \art{250} resuelve el caso del acreedor que es a la vez deudor de la masa: las sumas que le
correspondan se aplican al pago de su deuda \destacar{aunque no estuviere vencida}, operando una
suerte de compensación concursal.

\subsection{Resolución de controversias: el incidente del artículo 131}

El \art{131} establece el mecanismo general para resolver las cuestiones que se susciten entre el
deudor, el liquidador y cualquier interesado sobre la administración de los bienes: se tramitan
en \destacar{audiencias verbales} a solicitud del interesado, que expone por escrito su petición
y los antecedentes que la sustentan. Es la vía incidental a la que remiten numerosas
disposiciones del procedimiento ---entre ellas las controversias sobre incautación de
remuneraciones (capítulo \ref{cap:liquidacion-simplificada})---.

\subsection{La continuación de las actividades económicas}

La liquidación no siempre implica el cese inmediato de la operación. El \art{231} regula la
\destacar{continuación de actividades económicas}, en dos modalidades:

\begin{description}[leftmargin=0pt,style=nextline,font=\sffamily\bfseries]
  \item[Provisional] La decide el \emph{liquidador} para aumentar el porcentaje de recuperación,
  ejecutar prestaciones pendientes beneficiosas para la masa, o \destacar{propender a la
  realización de los activos como unidad económica}. El \art{232} le impone informar al tribunal
  y a la \superir{}, al día siguiente, las razones, los bienes adscritos y la fecha de inicio.

  \item[Definitiva] La acuerda la \emph{junta de acreedores}, con las mayorías y por el plazo que
  ésta determine.\verificar{Confirmar el artículo y el quórum de la continuación definitiva.}
\end{description}

\begin{foco}[Continuar operando para vender mejor]
La continuación no busca salvar la empresa ---eso es la reorganización--- sino \emph{vender más
caro}. Una unidad económica en funcionamiento vale más que sus activos apagados y dispersos. La
continuación provisional es la herramienta del liquidador para preservar ese mayor valor
mientras prepara la venta como unidad económica.
\end{foco}

\subsection{La realización ordinaria de los bienes}

El \art{207} entrega a la \destacar{junta de acreedores} la determinación de la forma, los
plazos y las condiciones de realización. El \art{208} enumera las fórmulas: venta al martillo
de muebles e inmuebles; remate en bolsa de los valores con presencia bursátil; u otra forma
distinta, incluida la venta como unidad económica.

\subsubsection{El silencio de los acreedores: la regla del artículo 210}

Aquí está la regla que conviene tener presente: si dentro de los \destacar{sesenta días}
contados desde la junta constitutiva la junta \emph{no acuerda} la forma de enajenación, los
bienes se enajenan \destacar{necesariamente conforme a las reglas de la realización sumaria o
simplificada} (arts. 203 y siguientes, capítulo \ref{cap:liquidacion-simplificada}).

\begin{foco}[La pasividad activa el modo simplificado]
El sistema no tolera la parálisis: si los acreedores no deciden en sesenta días, la ley decide
por ellos y aplica la realización sumaria. Para el liquidador es una salvaguarda ---no queda
atrapado sin forma de vender---; para el acreedor que prefiere otra modalidad, es un plazo fatal:
no pronunciarse equivale a aceptar la venta simplificada.
\end{foco}

\subsubsection{El martillero concursal}

El \art{213} reserva la facultad de rematar a los martilleros incluidos en la
\destacar{nómina que lleva la \superir{}}. El \art{214} dispone que el martillero se designa de
una \destacar{terna que propone el liquidador} con martilleros de esa nómina, y que las
condiciones de venta constan en las bases aprobadas por la junta. Concluido el remate, el
\art{216} obliga al martillero a rendir cuenta detallada ante la \superir{} \destacar{dentro de
quinto día} y a publicarla en el \boletin{}.

\subsubsection{Plazos y deber de información}

Rigen los plazos del \art{209} ---cuatro meses para muebles, siete para inmuebles---. El
\art{211} impone al liquidador comunicar a la \superir{}, con al menos quince días de
anticipación, la imposibilidad de cumplirlos, bajo sanción de \destacar{falta grave}.

\subsection{La venta como unidad económica}

El \art{217} permite a la junta acordar la \destacar{venta de un conjunto de bienes como unidad
económica}, cualquiera sea su naturaleza. Su efecto más potente lo fija el \art{218}: acordada
la enajenación como unidad, se \destacar{suspende el derecho de los acreedores hipotecarios,
prendarios y retencionarios} a realizar separadamente los bienes comprendidos en la unidad.

El \art{221} regula el cierre: la venta consta en escritura pública aprobada por el tribunal, el
cual ordena el \destacar{alzamiento y cancelación de todos los gravámenes y prohibiciones} sobre
los bienes que integran la unidad económica.

\begin{foco}[La unidad económica somete a los acreedores con garantía real]
Es una de las pocas instituciones que doblega la acción individual del acreedor hipotecario o
prendario: acordada la venta como unidad, no puede rematar por su cuenta el bien gravado. A
cambio, su preferencia se traslada al precio. Para el acreedor garantizado, esto convierte la
discusión en una sobre \emph{cuánto del precio le corresponde}, no sobre si ejecuta o no.
\end{foco}

\subsection{El contrato de \emph{leasing} en la liquidación}

El \art{225} resuelve la suerte del \destacar{arrendamiento con opción de compra}: la resolución
de liquidación \destacar{no es causal de terminación inmediata}. La junta constitutiva debe
pronunciarse entre continuar el contrato en los términos pactados, ejercer anticipadamente la
opción, o acordar con el arrendador una fórmula de realización que incluya los bienes (\art{227},
con quórum calificado).

Conforme al \art{226}, el arrendador puede verificar las cuotas devengadas e impagas hasta la
resolución de liquidación; las que se devenguen entre esa fecha y la junta constitutiva son
\destacar{de cargo de la masa}.

\begin{alerta}[El leasing no se extingue con la liquidación]
Un error frecuente es asumir que la liquidación termina automáticamente el \emph{leasing} y que
el bien puede incautarse como propio del deudor. No es así: el bien es del arrendador y el
contrato subsiste hasta que la junta decida. Incautar el bien como propio genera conflicto con
el arrendador y responsabilidad para el liquidador.
\end{alerta}

\subsection{El régimen recursivo restrictivo}

El \art{4} consagra la regla general: las resoluciones sólo son susceptibles de recurso
cuando la ley expresamente lo concede, y la apelación, cuando procede, se otorga en el solo
efecto devolutivo, salvo excepción legal.

La \destacar{Resolución de Liquidación} tiene un régimen recursivo propio y restringido,
coherente con la finalidad de celeridad. La \destacar{resolución de término} es apelable en
el solo efecto devolutivo, conservando entretanto el deudor la libre administración de sus
bienes (\art{256}).

\subsection{El rescate de la liquidación: término por acuerdo de reorganización}

La liquidación no es necesariamente un camino sin retorno. El \art{257} permite al deudor,
\destacar{durante el procedimiento de liquidación} y una vez notificada la nómina de créditos
reconocidos, \destacar{acompañar una propuesta de acuerdo de reorganización judicial}, a la que
se aplican las normas del Capítulo III en lo procedente.

El quórum del \art{258} es más exigente que el de la reorganización ordinaria: se requiere el
consentimiento del deudor y el voto conforme de \destacar{dos tercios de los acreedores
presentes que representen las tres cuartas partes del pasivo} con derecho a voto de la clase ---
frente a los dos tercios de pasivo de la reorganización del \art{79}---. Aprobado el acuerdo, el
\art{259} dispone que, en la misma resolución que lo tiene por aprobado, el tribunal declara el
\destacar{término legal del procedimiento de liquidación}.

\begin{foco}[Se puede revertir una liquidación en curso]
Aún abierta la liquidación, el deudor puede proponer un acuerdo que la deje sin efecto. El
precio es un quórum mayor ---tres cuartas partes del pasivo, no dos tercios---, lógico porque se
está revirtiendo una liquidación ya declarada. Para el deudor cuya empresa conserva valor de
marcha pese a la apertura, es la última oportunidad de salvataje; para el acreedor, una
alternativa que puede rendir más que el reparto liquidatorio.
\end{foco}

%==============================================================================
\bloquePractico
%==============================================================================

\subsection{Iter de la liquidación forzosa}

\begin{flujograma}{Tramitación de la liquidación forzosa}{cap07-iter}
  \node[hito] (demanda) {DEMANDA DE LIQUIDACIÓN FORZOSA (\art{117})\\[2pt]
    \normalfont Causal + consignación de 100 UF (\art{118})};
  \node[nodo, below=of demanda] (audiencia)
    {AUDIENCIA INICIAL (\art{120})\\[2pt]
     \normalfont El deudor individualiza sus tres mayores acreedores};

  \coordinate (split) at ($(audiencia.south)+(0,-6mm)$);
  \node[favorable, text width=38mm, below left=12mm and 2mm of audiencia.south] (enerva)
    {CONSIGNA, SE ALLANA O SE ACOGE A REORGANIZACIÓN};
  \node[nodo, text width=38mm, below right=12mm and 2mm of audiencia.south] (opone)
    {SE OPONE\\[2pt]\normalfont Juicio de oposición: prueba};

  \node[favorable, text width=38mm, below=9mm of enerva] (evita)
    {SE EVITA LA LIQUIDACIÓN\\[2pt]\normalfont O se reconduce a reorganización};
  \node[adverso, text width=38mm, below=9mm of opone] (resol)
    {RESOLUCIÓN DE LIQUIDACIÓN (\art{129})\\[2pt]
     \normalfont Desasimiento e incautación};

  \node[adverso, text width=38mm, below=9mm of resol] (realiza)
    {VERIFICACIÓN (30 días), JUNTA, REALIZACIÓN Y REPARTO};

  \draw[flecha] (demanda) -- (audiencia);
  \draw[flecha] (audiencia.south) -- (split) -| (enerva.north);
  \draw[flecha] (audiencia.south) -- (split) -| (opone.north);
  \draw[flecha] (enerva) -- (evita);
  \draw[flecha] (opone) -- (resol);
  \draw[flecha] (resol) -- (realiza);
\end{flujograma}

\subsection{Calendario de plazos}

\begin{table}[htbp]
\centering
\small
\caption{Plazos de la liquidación ordinaria}
\label{tab:plazos-liquidacion}
\begin{tabularx}{\textwidth}{@{}>{\raggedright\arraybackslash}p{0.38\textwidth} p{0.16\textwidth} X@{}}
\toprule
\textbf{Actuación} & \textbf{Plazo} & \textbf{Norma} \\
\midrule
Tramitación de la liquidación voluntaria & 3.\textsuperscript{er} día & \art{116} \\
\addlinespace
Presentar bienes suficientes (causal 2) & 4 días desde el requerimiento & \art{117} \Nro 2 \\
\addlinespace
Consignación a la orden del tribunal & 100 UF & \art{118} \Nro 2 \\
\addlinespace
Verificación ordinaria de créditos & 30 días desde la notificación & \art{170} \\
\addlinespace
Publicación de la nómina por el liquidador & 2 días de vencida la verificación & \art{170} \\
\bottomrule
\end{tabularx}
\end{table}

\subsection{Defensa del deudor en el juicio de oposición}

Remisión al capítulo \ref{cap:tram-liq-ordinaria}, que desarrolla la preparación de la
audiencia inicial, la decisión entre las cuatro opciones del \art{120} y la prueba de
solvencia.

%==============================================================================
\bloqueJurisprudencial
\label{sec:cap07-const}
%==============================================================================

\subsection{Constitucionalidad de la restricción del recurso de apelación}

El \art{4} \Nro 2 limita el recurso de apelación únicamente a las resoluciones que la propia
ley señala expresamente. Esta restricción ha sido cuestionada de forma reiterada ante el
\tconst{} por la vía de la inaplicabilidad por inconstitucionalidad, bajo el argumento de que
priva al deudor del control de un tribunal superior sobre decisiones críticas ---
señaladamente la declaración de liquidación forzosa--- y vulnera con ello el derecho a un
procedimiento racional y justo y la igualdad ante la ley.

\subsubsection{El caso de la falta de emplazamiento: TC Rol 11.421-2021}

El pronunciamiento más ilustrativo es el \rol{11.421-2021} \parencite{juris-tc-11421-2021}. La
requirente, Matilde Carolina Soto Rubilar, enfrentaba una liquidación forzosa promovida por
Comercial Soquimich S.A. ante el Segundo Juzgado de Letras de Chillán. El vicio era grave: la
demanda consignó un \destacar{domicilio inexistente}, y la notificación se practicó por la vía
del artículo 44 del \cpc{} en rebeldía absoluta de la deudora, que se enteró indirectamente de la
liquidación ya declarada. Promovido un incidente de nulidad por falta de emplazamiento y
desestimado, la aplicación del \art{4} \Nro 2 \destacar{le cerraba toda apelación} ante la CA
de Chillán.

El \tconst{} \destacar{acogió la inaplicabilidad}: privar de forma absoluta del recurso frente a
una alegación tan sustancial como la falta de emplazamiento configura una indefensión
incompatible con la Carta Fundamental. La celeridad concursal ---sostuvo--- no justifica
sacrificar la defensa de un deudor que nunca fue válidamente llamado a juicio. La misma
tensión se examinó en los \rol{12.539-2021} y \rol{12.527-2021}, de 3 de noviembre de 2022
\parencite{juris-tc-12539-2021, juris-tc-12527-2021}, y en el \rol{9.097-2020-INA}, de 21 de
abril de 2022 \parencite{juris-tc-9097-2020}.

\begin{foco}[La rigidez del artículo 4 cede ante el vicio de emplazamiento]
La doctrina del TC ha forzado una morigeración: las Cortes de Apelaciones reconocen que la
restricción recursiva \emph{cede} cuando está comprometida la regularidad misma de la relación
procesal ---la notificación, el emplazamiento---. Para el deudor mal notificado, la vía no es
sólo el incidente de nulidad: es la inaplicabilidad del \art{4} \Nro 2 ante el TC, que abre el
recurso que la ley concursal cerró.
\end{foco}

\subsubsection{El caso de la incautación de remuneraciones: TC Rol 10.957-2021}

Un supuesto conexo se planteó en el \rol{10.957-2021} \parencite{juris-tc-10957-2021}, promovido
por Raúl Juan Carlos González Guerra en su liquidación voluntaria ante el Juzgado de Letras de
Tocopilla. La liquidadora incautó \destacar{la totalidad de sus remuneraciones} ---\$7.671.442,
correspondientes a bonos extraordinarios de una negociación sindical--- de la cuenta corriente,
sin fijar bases de cálculo ni respetar los mínimos de subsistencia. El \art{131} inciso final
sólo admitía reposición en la misma audiencia, vedando la apelación; combinado con el \art{276},
dejaba un acto de disposición patrimonial de esa envergadura al arbitrio único del juez de la
instancia. El requirente alegó la infracción de su derecho de propiedad y a la subsistencia
digna.

\begin{foco}[El eje del debate: celeridad frente a doble instancia]
El régimen recursivo restrictivo sirve a la celeridad ---valor central del concurso--- pero a
costa de la doble instancia. El Tribunal Constitucional no ha invalidado el sistema en
abstracto: examina caso a caso si, en la gestión concreta, la denegación del recurso deja al
justiciable sin defensa idónea frente a una resolución gravosa. La estrategia de impugnación
debe, por tanto, construirse sobre el vicio específico de la resolución concreta ---la falta de
emplazamiento, la incautación del sustento vital--- y no sobre la inconstitucionalidad general
del artículo 4.
\end{foco}

\subsection{El desasimiento y la nulidad de los actos posteriores}

\pendiente{Sistematizar la jurisprudencia sobre la nulidad de los actos del deudor
posteriores a la resolución (\art{130} \Nro 1), su naturaleza y los efectos frente a
terceros de buena fe.}

\subsection{Responsabilidad del liquidador en la realización}

\pendiente{Remitir al capítulo \ref{cap:institucionalidad} y sistematizar los fallos sobre
responsabilidad civil del liquidador por devaluación de bienes no realizados oportunamente.}
